\documentclass[12pt]{article}
\usepackage[T1]{fontenc}
\usepackage[utf8]{inputenc}
\usepackage{lmodern}
\usepackage{textcomp}
\usepackage{enumitem}
\usepackage{geometry}
\geometry{margin=1in}
\begin{document}
\begin{center}
Answer Key - Economics - Undergraduate Level Assessment 12 \
\small (Answers are ordered exactly as in the question paper.)
\end{center}

\section*{Multiple Choice}
\begin{enumerate}[label=\arabic*.]
\item \textbf{Answer:} A
\\ \textit{Options:} A) Decrease of \$500 billion in NNP; B) Increase of \$500 billion in NNP; C) Decrease of \$100 billion in NNP; D) No change in NNP because taxes cancel out expenditures; E) Increase of \$200 billion in NNP
\\ \textit{Note:} The correct answer is based on the multiplier effect, where a tax increase of $100 billion reduces disposable income and leads to a total decrease in Net National Product (NNP) of $500 billion due to the marginal propensity to consume (MPC) of 0.8, resulting in a multiplier of 5 (1/(1-MPC)).
\item \textbf{Answer:} A
\\ \textit{Options:} A) Interest rates increase and the dollar depreciates.; B) Interest rates remain the same and the dollar depreciates.; C) Interest rates decrease and the dollar appreciates.; D) Interest rates and the value of the dollar both remain the same.; E) Interest rates increase but the value of the dollar remains the same.
\\ \textit{Note:} Contractionary fiscal policy, which involves reducing government spending or increasing taxes, leads to higher interest rates as the government borrows less, attracting investment and making the dollar more expensive, but simultaneously can lead to a depreciation of the dollar due to decreased economic growth expectations and reduced demand for U.S. exports.
\item \textbf{Answer:} C
\\ \textit{Options:} A) Do not produce if the TFC is not covered by revenue.; B) Do not produce if the price is lower than AVC.; C) Compare AVC to MR.; D) Compare AFC to MR.; E) Shut down if the firm is not making a profit.
\\ \textit{Note:} The correct answer is based on the principle that a competitive firm should continue operating in the short run as long as its marginal revenue (MR) covers its average variable costs (AVC), since this allows the firm to minimize losses rather than shutting down and incurring the full loss of fixed costs.
\item \textbf{Answer:} C
\\ \textit{Options:} A) The price of coffee machines rises.; B) The wage of coffee plantation workers falls.; C) The price of tea rises.; D) A decrease in the popularity of energy drinks.; E) Technology in the harvesting of coffee beans improves.
\\ \textit{Note:} The correct answer is based on the concept of substitute goods, as an increase in the price of tea, a substitute for coffee, would lead consumers to shift their preferences toward coffee, thereby increasing its demand.
\item \textbf{Answer:} A
\\ \textit{Options:} A) They signify shifts in currency exchange rates; B) They are indicative of immediate changes in consumer confidence; C) They determine the exact size of the workforce; D) They indicate economic stability; E) They show population change
\\ \textit{Note:} Slight differences in growth rate percentages are crucial because they can indicate changes in economic conditions that impact currency exchange rates, affecting international trade and investment dynamics.
\end{enumerate}

\section*{True / False}
\begin{enumerate}[label=\arabic*.]
\setcounter{enumi}{5}
\item \textbf{Answer:} False
\\ \textit{Options:} A) True; B) False
\\ \textit{Note:} The elasticity of supply is typically greater when producers have more capital to invest because greater capital allows for increased production capacity and flexibility in response to price changes.
\item \textbf{Answer:} True
\\ \textit{Options:} A) True; B) False
\\ \textit{Note:} When the U.S. dollar appreciates, it increases the purchasing power of U.S. residents abroad, making foreign goods and services, including vacations, cheaper in comparison to domestic prices.
\item \textbf{Answer:} False
\\ \textit{Options:} A) True; B) False
\\ \textit{Note:} Wage earners may not necessarily gain from anticipated inflation because, while wages can rise, they often do not keep pace with the rate of inflation, leading to a decrease in real purchasing power.
\item \textbf{Answer:} False
\\ \textit{Options:} A) True; B) False
\\ \textit{Note:} Models in economics can simplify complex realities and rely on assumptions that may not hold true in practice, potentially leading to inaccurate data interpretations and suboptimal decision-making outcomes for policymakers.
\item \textbf{Answer:} False
\\ \textit{Options:} A) True; B) False
\\ \textit{Note:} The statement is false because while international trade and competition can influence agricultural policies, factors such as environmental sustainability, food security, and rural development are also significant considerations for advocating special policies for farmers.
\end{enumerate}

\section*{Long Form Response}
\begin{enumerate}[label=\arabic*.]
\setcounter{enumi}{10}
\item \textbf{Answer:} The relationship between average product (AP) and marginal product (MP) of labor is crucial for understanding production efficiency. When the marginal product of labor exceeds the average product, it pulls the average up, leading to an increase in AP. Conversely, when MP falls below AP, it causes the average product to decline. This dynamic indicates that efficient production occurs when the marginal contributions of additional labor are maximized, aligning with or exceeding the average productivity level. Therefore, monitoring changes in MP is essential for firms aiming to optimize their labor input and overall production efficiency.
\\ \textit{Note:} correct because it accurately describes how the marginal product influences the average product by indicating that an increase in marginal product raises the average product, while a decrease in marginal product lowers it, thus illustrating the importance of maximizing labor contributions for production efficiency.
\item \textbf{Answer:} In a monopolistically competitive firm's profit-maximizing equilibrium, the inefficiency arises primarily from the fact that the price charged by the firm exceeds its marginal cost. This scenario leads to a situation where the quantity of output produced is less than what would be socially optimal, as the firm restricts production to maximize profits. When marginal cost is below price, it indicates that the firm could increase output without incurring additional costs that would exceed the price consumers are willing to pay, thereby suggesting that society is not receiving enough goods. This underproduction results in a deadweight loss, as potential gains from trade are not fully realized, ultimately diminishing societal welfare. Thus, the divergence between price and marginal cost in monopolistic competition leads to an overall inefficiency in resource allocation and output.
\\ \textit{Note:} correct because it accurately identifies that in monopolistic competition, firms maximize profits by setting prices above marginal costs, resulting in reduced output and a loss of societal welfare compared to the socially optimal level of production.
\end{enumerate}

\vfill
\noindent\textit{End of Answer Key}
\end{document}